\chapter{Inventar de endpoint-uri API}

\section{Observații generale}

Această anexă sintetizează endpoint-urile principale ale aplicației. Nu urmăresc să refac documentația Swagger în întregime, ci să ofer o imagine clară asupra suprafeței API-ului și a modului în care aceasta este împărțită pe domenii funcționale.

\section{Endpoint-uri pentru cont și autentificare}

\begin{longtable}{>{\RaggedRight\arraybackslash}p{2cm}>{\RaggedRight\arraybackslash}p{5.7cm}>{\RaggedRight\arraybackslash}p{6.2cm}}
  \caption{Endpoint-uri principale din zona de autentificare} \\
  \toprule
  \textbf{Metodă} & \textbf{Rută} & \textbf{Rol} \\
  \midrule
  \endfirsthead
  \toprule
  \textbf{Metodă} & \textbf{Rută} & \textbf{Rol} \\
  \midrule
  \endhead
  POST & `/api/account/register` & creare cont nou \\
  POST & `/api/account/login` & autentificare clasică \\
  POST & `/api/account/oauth2/github` & autentificare prin GitHub \\
  POST & `/api/account/oauth2/google` & autentificare prin Google \\
  POST & `/api/account/two-factor/verify-login` & finalizarea login-ului cu 2FA \\
  POST & `/api/account/oauth2/github/link` & legare cont GitHub la utilizatorul curent \\
  POST & `/api/account/oauth2/google/link` & legare cont Google la utilizatorul curent \\
  POST & `/api/account/forgot-password` & inițiere resetare parolă \\
  POST & `/api/account/reset-password` & setarea noii parole prin token \\
  POST & `/api/account/set-password` & adăugarea unei parole pentru cont OAuth \\
  POST & `/api/account/change-password` & schimbarea parolei \\
  POST & `/api/account/two-factor/enable/request` & trimiterea codului pentru activare 2FA \\
  POST & `/api/account/two-factor/enable/confirm` & confirmarea activării 2FA \\
  POST & `/api/account/two-factor/disable/request` & trimiterea codului pentru dezactivare 2FA \\
  POST & `/api/account/two-factor/disable/confirm` & confirmarea dezactivării 2FA \\
  POST & `/api/account/refresh` & reînnoirea token-ului de acces \\
  POST & `/api/account/logout` & încheierea sesiunii \\
  \bottomrule
\end{longtable}

\section{Endpoint-uri pentru proiecte}

\begin{longtable}{>{\RaggedRight\arraybackslash}p{2cm}>{\RaggedRight\arraybackslash}p{5.7cm}>{\RaggedRight\arraybackslash}p{6.2cm}}
  \caption{Endpoint-uri principale pentru proiecte} \\
  \toprule
  \textbf{Metodă} & \textbf{Rută} & \textbf{Rol} \\
  \midrule
  \endfirsthead
  \toprule
  \textbf{Metodă} & \textbf{Rută} & \textbf{Rol} \\
  \midrule
  \endhead
  GET & `/api/projects` & lista proiectelor utilizatorului curent \\
  GET & `/api/projects/{id}` & proiect după identificator \\
  GET & `/api/projects/by-slug/{slug}` & proiect după slug \\
  POST & `/api/projects` & creare proiect \\
  PUT & `/api/projects/{id}` & actualizare metadate proiect \\
  DELETE & `/api/projects/{id}` & ștergere proiect \\
  POST & `/api/projects/{id}/assets/images` & upload imagine pentru canvas \\
  GET & `/api/projects/{id}/design` & încărcare design salvat \\
  PUT & `/api/projects/{id}/design` & salvare design \\
  POST & `/api/projects/{id}/flush` & salvare combinată design + thumbnail \\
  PUT & `/api/projects/{id}/thumbnail` & actualizare thumbnail \\
  GET & `/api/projects/user/{userId}` & proiecte publice ale unui utilizator \\
  POST & `/api/projects/{id}/star` & bookmark / star \\
  DELETE & `/api/projects/{id}/star` & eliminare bookmark \\
  POST & `/api/projects/{id}/view` & înregistrare vizualizare \\
  POST & `/api/projects/{id}/like` & apreciere proiect \\
  DELETE & `/api/projects/{id}/like` & eliminare apreciere \\
  POST & `/api/projects/{id}/fork` & creare fork după proiect public \\
  \bottomrule
\end{longtable}

\section{Endpoint-uri pentru utilizatori și profiluri}

\begin{longtable}{>{\RaggedRight\arraybackslash}p{2cm}>{\RaggedRight\arraybackslash}p{5.7cm}>{\RaggedRight\arraybackslash}p{6.2cm}}
  \caption{Endpoint-uri principale pentru utilizatori} \\
  \toprule
  \textbf{Metodă} & \textbf{Rută} & \textbf{Rol} \\
  \midrule
  \endfirsthead
  \toprule
  \textbf{Metodă} & \textbf{Rută} & \textbf{Rol} \\
  \midrule
  \endhead
  GET & `/api/users/me` & profilul utilizatorului curent \\
  PUT & `/api/users/me` & actualizare profil propriu \\
  POST & `/api/users/me/profile-image` & upload imagine de profil \\
  DELETE & `/api/users/me` & ștergere cont propriu \\
  DELETE & `/api/users/me/linked-accounts/{provider}` & deconectare provider OAuth \\
  GET & `/api/users/search?q=...` & căutare de utilizatori \\
  GET & `/api/users/{username}` & profil public după username \\
  POST & `/api/users/{username}/follow` & urmărire utilizator \\
  DELETE & `/api/users/{username}/follow` & oprire urmărire \\
  GET & `/api/users/{username}/followers` & lista urmăritorilor \\
  GET & `/api/users/{username}/following` & lista utilizatorilor urmăriți \\
  GET & `/api/users/me/stars` & proiectele salvate de utilizatorul curent \\
  \bottomrule
\end{longtable}

\section{Endpoint-uri pentru explore și recomandări}

\begin{longtable}{>{\RaggedRight\arraybackslash}p{2cm}>{\RaggedRight\arraybackslash}p{5.7cm}>{\RaggedRight\arraybackslash}p{6.2cm}}
  \caption{Endpoint-uri din zona explore} \\
  \toprule
  \textbf{Metodă} & \textbf{Rută} & \textbf{Rol} \\
  \midrule
  \endfirsthead
  \toprule
  \textbf{Metodă} & \textbf{Rută} & \textbf{Rol} \\
  \midrule
  \endhead
  GET & `/api/explore/trending` & proiecte populare \\
  GET & `/api/explore/recommended` & recomandări de proiecte, eventual personalizate \\
  GET & `/api/explore/people` & sugestii de utilizatori \\
  \bottomrule
\end{longtable}

\section{Endpoint-uri pentru AI și conversie}

\begin{longtable}{>{\RaggedRight\arraybackslash}p{2cm}>{\RaggedRight\arraybackslash}p{5.7cm}>{\RaggedRight\arraybackslash}p{6.2cm}}
  \caption{Endpoint-uri pentru AI și converter} \\
  \toprule
  \textbf{Metodă} & \textbf{Rută} & \textbf{Rol} \\
  \midrule
  \endfirsthead
  \toprule
  \textbf{Metodă} & \textbf{Rută} & \textbf{Rol} \\
  \midrule
  \endhead
  POST & `/api/ai/design` & generare design într-o singură cerere \\
  POST & `/api/ai/design/stream` & generare design în regim SSE \\
  POST & `/api/ai/design/pipeline` & rulare pipeline AI în trei faze \\
  POST & `/api/ai/design/pipeline/stream` & rulare pipeline AI cu streaming \\
  POST & `/api/converter/generate` & export pentru o pagină sau set responsive \\
  POST & `/api/converter/validate` & validare structură IR \\
  POST & `/api/converter/generate-files` & generare fișiere pentru export multi-page \\
  \bottomrule
\end{longtable}

\noindent\textbf{Observație finală.}

Din inventarul de mai sus se vede că API-ul nu este limitat la o singură zonă funcțională. El trebuie să deservească atât fluxuri de produs clasice, cât și funcționalități tehnice speciale, cum sunt streaming-ul AI și exportul multi-framework. Tocmai această diversitate justifică structura pe straturi din backend și separarea convertorului în proiect propriu.
